\section{Preliminaries}
\subsection{Notation and Terminology}\label{sc:notation}
Throughout this paper, neighbors of a node $v$ are denoted by $N(v)$. We denote by $d_G(u,v)$ the distance between $u$ and $v$ in the graph $G$. If $G$ is clear from the context, we use the notation $d(u,v)$. We denote $W$ for the ratio between the maximum and minimum weight of the graph.  

\paragraph{Dynamic graph algorithms.}
We are interested in designing dynamic graph algorithms for shortest path problems, which allow both for update and query operations. In particular, we look at \emph{decremental} algorithms, which allow for weight increases and edge deletions. This is opposed by \emph{fully dynamic} algorithms, which also allow for weight decreases and edge insertions. To be concrete, we allow for the following three operations:
\begin{itemize}
    \item \Delete{$u,v$}: delete the edge $\{u,v\}$ from the graph;
    \item \Update{$u,v,\Delta$}: increase the weight of an edge $\{u,v\}$ to $\Delta$, i.e., $w(u,v)\leftarrow \Delta $;
    \item \Distance{$u,v$}: return a distance approximation $\hat d_G(u,v)$ between the nodes $u$ and $v$ in the current graph $G$.
\end{itemize}

We assume that we have an \emph{oblivious adversary}, i.e., that the sequence of updates is fixed from the start. This is as opposed to the \emph{adaptive adversary}, which can use the random choices made by the algorithm to determine the next update. 

Edge weights can be positive integers bounded by $W$, for some parameter $W\in \mathbb N$. We say that $\Distance(u,v)$ returns \emph{$(\alpha,\beta)$-approximate distances}, denoted by $\hat d_G$, if $d_G(u,v) \leq \hat d_G(u,v) \leq \alpha d_G(u,v)+\beta$ for all $u,v\in V$. Moreover we denote $(\alpha,0)$-approximate by $\alpha$-approximate. 

In all our algorithms, we maintain a distance oracle with constant \emph{query time}, i.e., the time needed to answer a distance query $\Distance{u,v}$.
Further, we give the \emph{total update time}, which is the time needed to process a sequence of up to $m$ edge deletions or weight increases, where $m$ is the number of edges before the first deletion. 

For any object, value, or data structure $x$ that changes over time, we use a superscript $x^{(t)}$ to specify its state at time $t$, i.e., $x^{(t)}$ is $x$ after processing $t$ updates. 



\subsection{Decremental Tools}\label{sc:decr_tools}
We start by describing several tools from previous work that we need to use in our new decremental algorithm.
\paragraph{Cluster and bunch maintenance.} 
Recall Thorup and Zwick~\cite{TZ2005} clusters as follows: Given a parameter $p\in [0,1]$, called the \emph{cluster sampling rate}, we let $A$ be a set of sampled nodes, where each node is sampled uniformly at random with probability $p$. Then we define the \emph{pivot} of a node $v$ to be the closest node in $A$, i.e., $p(v)=\argmin_{s\in A} d(s,v)$. 

The following lemma is implicitly proven in \cite{LN2020}, but for completeness we sketch a proof in Appendix~\ref{ap:bunches_and_clusters}.

\begin{restatable}{lemma}{lmbunch} \label{lm:approx_bunches_on_G}
Given a decremental graph $G=(V,E,w)$, and parameters $0< \epsilon <1$, $0<\rho<1/2$ and the sampling probability $0< p <1$ we can maintain for each $v\in V$ an approximate pivot $\tilde p(v)$, such that $d_G(v,p(v)) \leq \dLNbunch(v,\tilde{p}(v)) \leq (1+\epsilon) d_G(v,p(v))$.
Moreover, we maintain approximate bunches $\tilde{B}(v)$ that contains
$\hat{B}(v):= \{ w \in V : (1+\epsilon)d_G(v,w) < d_G(v, p(v)) \}$ and we have monotone (non-decreasing over updates) distance estimates $\dLNbunch$ such that for any $u \in \tilde{B}(v)$ we have $d_G(u,v) \leq \dLNbunch(u,v) \leq (1+\epsilon) d_G(u,v)$, and with high probability, we have $|\tilde{B}(v)| = O(\frac{\log n\log nW}{p})$. The algorithm has total update time $\tilde O(\tfrac{mn^\rho}{p}) \cdot O(\frac{1}{\rho \epsilon})^{1+2/\rho}$.
\end{restatable}

In our algorithms for balancing out different terms in the running time we set the parameter~$\rho$ depending on the graph density so that $m= n^{1+\rho}$. Hence for graphs with $m= n^{1+\Omega(1)}$ edges, we can set $\rho$ to be a constant and for sparser graphs we can set $\rho = \frac{\sqrt{\log n}}{\log \log n}$. For simplicity of our statements in the rest of this section we assume $0< \epsilon <1$ is a constant, and thus we hide $(1/\epsilon)^c$ factors in $\tilde{O}$ when $c$ is constant, or in $n^{o(1)}$ when $c=\omega(1)$.



\paragraph{Multi-source shortest paths.}
We will use the following lemma (from \cite{LN2020}) repeatedly to maintain $(1+\epsilon)$ approximate shortest paths from a given set of sources. 

\begin{lemma}[\cite{LN2020}]\label{lm:decr_MSSP}
    There is a data structure which given a weighted undirected graph $G = (V, E)$ explicitly maintains $(1+\epsilon)$-approximate distances from a set of sources $S$ in $G$ under edge deletions, where $0<\epsilon< 1/2$ is a constant. After each update, the algorithm returns all nodes for which the distance estimate has changed. 
    During the execution of the algorithm, nodes can be added to $S$. We write $s$ for the final size of $S$. 
    Assuming that $|E| = n^{1+\Omega(1)}$ and $s = n^{\Omega(1)}$, the total update time is w.h.p.~$\tilde O(sm)$. If $|E| = n^{1+o(1)}$ or $s= n^{o(1)}$, then the total update time is $\tilde O(smn^{o(1)})$. 
\end{lemma}

